% Potenz- und Exponentialfunktionen, Wachstum und Zerfall · Wert und Schwellenwert berechnen · Prüfungs-Fokus MSA – bank/potenz-exponentialfunktionen/e3.jsonl (zusammenbau v0.6, Prüfungs-Fokus)
\einheitenkopf[e3]{Einheit 3 · Exponentialfunktion aufstellen und auswerten}
\zweigzeile{ab Kl. 10 · P10 ×7}
\setcounter{aufgabe}{0}

% Anlauf (ohne Prüfkennung)
\begin{aufgabe}{Wert und Schwellenwert berechnen – Anlauf}
\begin{teile}
\teil Berechne $1{,}04^5$ und runde auf vier Stellen. ≈ \leerfeld
\teil Ein Bestand von $400$ wächst jedes Jahr mit dem Faktor $1{,}1$. Berechne die Werte nach $1$, $2$ und $3$ Jahren Schritt für Schritt. \leerfeld , \leerfeld , \leerfeld
\teil Eine Gemeinde hatte $2019$ $2\,500$ Einwohner; die Zahl wächst jährlich um $6\,\%$. Wie viele Einwohner sind es $2025$? Runde auf ganze Personen. t = \leerfeld , ≈ \leerfeld
\end{teile}
\end{aufgabe}

\begin{aufgabe}{Wert und Schwellenwert berechnen – Prüfungsaufgaben}
\begin{teile}
\steil Die Zahl der Hefezellen in einer Probe wächst exponentiell mit $N(t) = 64 \cdot 1{,}25^t$ ($t$ in Stunden). Aussage: „Nach $30$ Stunden sind es mehr als $50\,000$ Zellen.“ Prüfe die Aussage rechnerisch. \hfill (P25*)
\steil Die Nutzerzahl einer App verdoppelt sich jeden Monat: $N(t) = 120 \cdot 2^t$ ($t$ in Monaten). Aussage: „Nach $20$ Monaten hat die App mehr als hundert Millionen Nutzer.“ Prüfe die Aussage rechnerisch. \hfill (P25*)
\steil Studie 1 beschreibt den Stromverbrauch einer Firma mit $y = 36 \cdot 1{,}04^x$ ($x$ Jahre ab $2021$, $y$ in Tausend Kilowattstunden). Studie 2 sagt von $2021$ bis $2033$ insgesamt etwa $50\,\%$ Steigerung voraus. Vergleiche beide Vorhersagen für $2033$. \hfill (P20*)
\steil Studie 1 beschreibt den Wasserverbrauch eines Ortes mit $y = 120 \cdot 1{,}02^x$ ($x$ Jahre ab $2023$, $y$ in Tausend Kubikmetern). Studie 2 erwartet bis $2043$ insgesamt etwa $40\,\%$ mehr. Vergleiche beide Vorhersagen für $2043$. \hfill (P20*)
\steil Von einem Kontrastmittel sind zu Beginn $8{,}00\,\mathrm{mg}$ im Körper; jede Stunde werden $20\,\%$ abgebaut. Wie viel ist nach $24$ Stunden noch vorhanden? \hfill (P16*) \\ ≈ \leerfeld[mg]
\steil Von einem Farbstoff sind zu Beginn $80\,\mathrm{mg}$ in einer Lösung; jeden Tag zerfallen $25\,\%$. Wie viel ist nach $24$ Tagen noch vorhanden? \hfill (P16*) \\ ≈ \leerfeld[mg]
\teil Der Wasserdruck in einer Leitung sinkt je Kilometer um $12\,\%$; am Anfang sind es $800\,\mathrm{hPa}$. An einer Stelle werden $480\,\mathrm{hPa}$ gemessen. Behauptung: „Die Stelle liegt etwa $4\,\mathrm{km}$ vom Anfang entfernt.“ Entscheide mit Rechnung. \hfill (P17)
\teil Die Helligkeit unter Wasser nimmt je Meter Tiefe um $15\,\%$ ab; an der Oberfläche sind es $1\,000$ Lux. Ein Taucher misst $444$ Lux. Behauptung: „Er ist etwa $5\,\mathrm{m}$ tief.“ Entscheide mit Rechnung. \hfill (P17)
\teil Der Umsatz einer Bäckerei betrug $2023$ $612\,000\,€$ und wächst jährlich um $7\,\%$. In welchem Jahr liegt der Umsatz erstmals über $800\,000\,€$, und wie hoch ist er dann? \hfill (P18) \\ Jahr \leerfeld , ≈ \leerfeld[€]
\teil Eine Gemeinde hatte $2022$ $8\,240$ Einwohner; die Zahl wächst jährlich um $3\,\%$. In welchem Jahr sind es erstmals mehr als $9\,500$, und wie viele sind es dann? \hfill (P18) \\ Jahr \leerfeld , ≈ \leerfeld
\end{teile}
\end{aufgabe}

\medskip
\uebersichtskasten{Exponentialfunktion: Anfangswert mal Wachstumsfaktor hoch Anzahl der Schritte – die Variable steht im Exponenten: N(t) = N₀ · qᵗ; Anfangswert 200 und Faktor 1,3 ergeben N(t) = 200 · 1,3ᵗ. \\ Wert berechnen: Zahl der Schritte für t einsetzen und die Potenz mit dem Taschenrechner rechnen: N(4) = 200 · 1,3⁴ ≈ 571. \\ Schritte zählen: Der Startzeitpunkt ist t = 0; von 2020 bis 2024 sind es vier Schritte, nicht 2024. \\ Schwellenwert: Schritt für Schritt mal q rechnen, bis der Wert die Schwelle überschreitet, und die Schritte zählen.}
